%*******************************************************
% Abstract in English
%*******************************************************
\pdfbookmark[0]{Abstract}{Abstract}

\begin{otherlanguage}{american}
    \chapter*{Abstract}

    When conflicting information admits several valid assessments, the
    relevant question is not only which arguments can be accepted, but
    also whether partial decisions can be combined independently while a
    shared context is held fixed. Abstract argumentation frameworks (AFs)
    model such situations through arguments and directed attacks.
    Conditional independence asks whether the assessments of two sets of
    arguments \(X\) and \(Y\), taken from labellings that agree on a
    conditioning set \(Z\), can be recombined in another valid labelling.
    Deciding this property exactly is computationally demanding because
    of its universal-existential structure. Moreover, existing
    graph-based criteria identify only some of the independences that
    actually hold. This thesis therefore addresses the absence of an
    exact decision procedure whose practical behavior has been
    investigated empirically.

    A uniform encoding using quantified Boolean formulas (QBFs) is
    developed for the \emph{complete}, \emph{stable}, and
    \emph{preferred} semantics. Boolean variables represent the labels
    \(\mathsf{in}\), \(\mathsf{out}\), and \(\mathsf{und}\) of each
    argument. Semantic encodings, equality constraints, and a mixing
    formula connect two universally considered source labellings with an
    existentially required recombination labelling. For the
    \emph{preferred} semantics, the maximality of the accepted arguments
    is additionally encoded. A correctness proof establishes that the
    resulting QBF is true if and only if the requested conditional
    independence holds. Consequently, all valid labellings do not have to
    be generated explicitly before the query can be decided.

    The encodings are translated into the \texttt{QCIR-G14} format and
    evaluated with QuAbS. The empirical study comprises \(9\,450\)
    queries on \(105\) AFs from the ICCMA'25 benchmark collection, using
    five singleton argument pairs, six conditioning levels, and the
    three considered semantics. Of the \(7\,164\) decided queries,
    \(6\,986\), or \(97.52\,\%\), are classified as independent, whereas
    \(178\), or \(2.48\,\%\), are classified as dependent. Among the
    decided queries, larger conditioning sets are associated with less
    observed dependence, as well as with generally longer solver
    runtimes and lower solution rates. The runtime effects are more
    pronounced for the \emph{complete} and \emph{stable} semantics than
    for the \emph{preferred} semantics.

    For the benchmark collection studied, the results show that conditional
    independence predominates among the decided queries under all three
    semantics and that larger conditioning sets are associated with less
    observed dependence and greater practical computational effort. The
    complexity analysis complements this picture: every independence query
    has a positive answer in acyclic and, more generally, even-cycle-free
    AFs, whereas the decision problem is
    \(\Pi_2^{\mathsf P}\)-complete under all three semantics in
    odd-cycle-free AFs. Together, the exact QBF encoding, the empirical
    evaluation, and the complexity results provide a unified foundation 
    for analyzing conditional independence in argumentation frameworks.
\end{otherlanguage}
